% File: main_aaai27_v43_candidate.tex
% Current AAAI-27 main manuscript, synchronized with the active supplement.
\documentclass[letterpaper]{article}
\usepackage[submission]{aaai2027}
\usepackage[hyphens]{url}
\usepackage{graphicx}
\urlstyle{rm}
\def\UrlFont{\rm}
\usepackage{natbib}
\usepackage{caption}
\frenchspacing

\usepackage{booktabs}
\usepackage{amsmath}
\usepackage{amssymb}

\pdfinfo{
/TemplateVersion (2027.1)
}

\setcounter{secnumdepth}{0}

\title{Evidence-Sealed Registration of Owner-Sampled DP-SGD Routes}

\author{
Anonymous Submission
}
\affiliations{
}

\begin{document}

\maketitle

\begin{abstract}
A completed training run does not support a reported owner-level privacy claim
when valid components from incompatible differentially private stochastic
gradient descent (DP-SGD) routes share one label.  We call
this joint predicate route-label registration consistency.  In private
artificial intelligence (AI) training over owner-attributed records, one route
must jointly own adjacency, sampler, sensitivity/noise geometry, accountant,
executor, transformation, and evidence.  An evidence-sealed judgment enforces
it.  Each route implementation (handler) registers a complete core of nominal
semantics, exact types and callables, a route-local registry, and seven
source-resolved obligations.  Before dispatch, registration recomputes this
core and requires exact report, schema, and digest equality.

We instantiate Bernoulli/add-remove (P), fixed-size/replace-one (F), and
per-epoch random allocation/add-remove (A).  Across three route pairs, weak callable/type and
source-bound checks accept, respectively, 378 and 42 of 378 prespecified
non-original handler splices; the complete-core seal accepts none, while all
six original endpoints pass.  On frozen P-lifecycle evidence, strict
contract-only parsing rejects
36/52 targeted mutations and 5/8 reproduced failures, versus 52/52 and 8/8 for
the full lifecycle.  Adding A through branch-free handler dispatch leaves six
frozen P/F contracts and public plans exactly unchanged, and all nine
dataset--route profiles compile.  The mechanisms and accountants are prior work.
These are bounded honest-source conformance results, not a privacy proof,
formal refinement, safety for arbitrary plugins, or cryptographic attestation.
\end{abstract}

\section{Introduction}

Differentially private stochastic gradient descent (DP-SGD) trains AI models by
clipping contributions, adding Gaussian noise, and accounting for repeated
access \citep{abadi2016deep,dwork2014foundations,mironov2017renyi}.  A
differential privacy (DP) statement is meaningful only after specifying
adjacency, the unit that may change; National Institute of Standards and
Technology (NIST) guidance treats the privacy unit as core to a DP guarantee
\citep{near2025nist}.  Sequential preprocessing often turns each person's
stream into many windows.  Training and utility evaluation may still complete
although example-level accounting does not support the reported owner claim.

Owner-level sampling and clipping are established
\citep{chua2024mind,ponomareva2023dpfy}.  Opacus supplies a core Poisson
training stack \citep{yousefpour2021opacus}; Wang and Birrell analyze
fixed-size sampling \citep{wang2019subsampled,birrell2024fixed}; random
allocation and Balanced Iteration Subsampling (BIS) supply $k$-out-of-$T$
participation and stronger accountants
\citep{shenfeld2025randomallocation,feldman2026efficientallocation,
dong2026lessrandom}.  Typed DP systems, implementation tests, and execution
attestation address broader or different questions
\citep{hay2020opendp,zhang2020dpcheck,kong2024dpauditorium,
abdolmaleki2026veridp}.  We study a private-AI integration boundary when these
route-specific components share one extensible training interface:

\begin{quote}
What must a route registration bind so that mixing components from two
registered owner-DP-SGD routes cannot retain either route label?
\end{quote}

We call this property \emph{route-label registration consistency}: one exact
route must jointly own its adjacency, owner sampler, sensitivity/noise geometry,
accountant, executor, mapping/preprocessor, and evidence chain.  We test three
prior-work endpoints: Bernoulli/add-remove (P), fixed-size/replace-one (F), and
per-epoch random allocation/add-remove (A).  A preserved historical failure
used an epoch-shuffle/drop-last loader but described it by a per-step
fixed-population Hypergeometric accounting model without established exact law
correspondence.  We therefore classify the exact
accounting claim as unsupported; this classification does not by itself
establish a privacy violation.

A route is not a sampler name.  We model it nominally, with an exact route
identity, as
\[
 \rho=(\mathord{\sim},\mathcal S,\Delta,\nu,\mathcal A,\mathcal E),
\]
where the entries are adjacency, owner sampler, sum sensitivity, Gaussian
scale, accountant, and executor.  Each registered route implementation is a
handler.  Our registration judgment makes a route-local passing report commit
to the complete core.  The generic dispatcher contains no route-name or
route-type branch: it validates the selected handler, invokes its exact
parser/compiler, and checks the returned type and postcompile integrity method.
We test extensibility by adding A, then checking exact preservation of the old
P/F objects and all prespecified nonendpoint splices among the three route pairs.

Exact route-specific parsing is necessary but insufficient.  In our ablation,
a generic map checking only callable and type compatibility accepts every one
of 378 pairwise mixed, non-original
(\emph{nonendpoint}) handler splices.  A source-bound pre-seal that also
resolves callable source files still accepts 42.  The surviving cases splice a
valid parser, compiler, or
accountant/cross-check pair while keeping individually valid source witnesses.
This ablation isolates the \emph{complete core} from another field-level check.
The contribution is the route-label predicate and its measured counterfactual,
not a one-route wrapper or a hashing technique.
SHA-256 is only the deterministic commitment used to enforce the predicate.
Software-supply-chain provenance can host this check, and reproducible builds
can strengthen its artifact chain; neither defines the DP predicate by itself.
We claim neither authenticated provenance nor independently reproducible
builds \citep{torresarias2019intoto,lamb2022reproducible}.

Two simpler explanations require separate controls.  First, perhaps strict
contract parsing already explains the lifecycle.  On 52 targeted P mutations
and eight reproduced historical failures fixed before stage attribution, a
strong contract-only boundary catches 36 and five; mapping, execution,
integrity, and artifact boundaries catch the remaining 16 and three.  Second,
perhaps multiple routes are only closed branches.  A matched P/F experiment
holds 38 scientific-configuration fields fixed and rejects all 48 prespecified
bidirectional obligation substitutions; branch-free addition of A then leaves
the six frozen P/F outputs exact.  These are finite controls, not completeness
or arbitrary-route guarantees.

Our contributions are:
\begin{itemize}
    \item a route-label registration-consistency criterion, enforced by a
    complete-core judgment that binds nominal semantics, exact types and
    callables, a route-local registry, seven source-resolved obligations, and
    one passing report before dispatch;
    \item a counterfactual ladder against simpler boundaries: contract-only
    parsing leaves 16/52 targeted and 3/8 historical cases to downstream
    lifecycle gates, while weak and source-bound handler boundaries accept 378
    and 42 splices and the complete-core seal accepts none; and
    \item bounded route-separation and extension tests: with 38 P/F scientific
    configuration fields fixed, all 48 obligation substitutions fail, while
    branch-free addition of A preserves six exact P/F contract/plan outputs,
    with nine compilations, 2,286/2,286 prespecified contract-splice rejections,
    and independent execution replays.
\end{itemize}

\paragraph{Claim and phase boundary.}
We propose no new sampler, accountant, amplification theorem, owner mechanism,
general DP compiler or type system, privacy tester, or attestation protocol.
In particular, prior random-allocation/BIS work owns Route A and tighter
accounting.  Our object is the finite three-route registration boundary.
Concurrent work maps frozen run artifacts to the privacy-unit wording they
authorize
\citep{anonymous2026audit}.  Our phase instead fixes and executes a registered
route before training and performs no post-hoc privacy-unit conversion.  The
evidence model is honest execution, not arbitrary-code proof.  The studies
reuse public rows and subject attribution from the Human Activity Recognition
(HAR) dataset hosted by the University of California, Irvine (UCI) Machine
Learning Repository, not exact mechanism plans, numerical results, or evidence
chains.

\section{Setting and Registered Routes}

\paragraph{Owners and generated examples.}
Let a dataset be $\mathcal D=(D_u)_{u\in U}$, where $D_u$ is the collection
owned by person, patient, subject, device, or sequence $u$.  Preparation
produces records $w$ and an attribution map $\alpha(w)$.  Every registered route
requires
$\alpha(w)=\{u\}$: every generated record has exactly one owner.  Multi-owner
records need another analysis \citep{ganesh2025multiattribution} and are
rejected.

A contribution-policy rule is fixed before training and acts separately
within each owner.  It may inspect that owner's records under fixed public caps
on retention and per-step selection; it must not couple one owner's vector to
another owner's raw data.  A preprocessing transform is fixed as public side
information and is never fit inside the executor on protected input.

\paragraph{Common clipped update.}
For participating owner $u$ at step $s$, the executor applies the owner-local
rule, computes a mean-loss gradient $g_{u,s}$, and forms
\[
 \bar g_{u,s}=g_{u,s}\min\left\{1,
       \frac{C}{\lVert g_{u,s}\rVert_2}\right\}.
\]
Let $\theta_s$ be the model parameters at step
$s\in\{0,\ldots,T-1\}$, $B_s$ the participating-owner set, $\eta$ the
learning rate, $C$ the clipping threshold, $\sigma$ the noise multiplier, and
$I$ the identity matrix.
All routes execute the fixed schedule
\begin{equation}
\theta_{s+1}=\theta_s-\frac{\eta}{D}
\left(\sum_{u\in B_s}\bar g_{u,s}+Z_s\right),\quad
Z_s\sim\mathcal N(0,\sigma^2C^2 I).
\label{eq:update}
\end{equation}
$D$ is fixed post-processing, never the realized sample size.

\paragraph{Route P: Bernoulli/add-remove.}
The public plan fixes $q,T,C,\sigma,D$.  Every present owner is independently
included with probability $q$.  Adding one owner changes the clipped sum by at
most $C$.  When $B_s=\varnothing$, the sum is zero but noise is still sampled
and the optimizer advances.

\paragraph{Route F: fixed-size/replace-one.}
The public plan fixes a population of $N$ owner positions and
$m,T,C,\sigma,D=m$.  At every step it uses simple random sampling without
replacement (SRSWOR), drawing a uniform size-$m$ subset; the sampled positions
are distinct and $m>0$.  Replacing one
owner changes the clipped sum by at most
$\lVert\bar g-\bar g'\rVert_2\leq 2C$.  Thus noise standard deviation
$\sigma C$ corresponds to normalized Gaussian noise $\sigma/2$ for its
SRSWOR accountant.  This route does not inherit Route P's adjacency or
accountant.

\paragraph{Route A: random allocation/add-remove.}
For each of $E$ epochs, every owner independently draws a uniform $k$-subset
of $t$ candidate steps and participates exactly at those steps, giving total
schedule length $T=Et$.  Thus one owner's
participation is coupled across steps, unlike P, while the number active at a
step is not fixed, unlike F.  Every step adds noise and advances the optimizer,
including an empty step.  The sum has add/remove sensitivity $C$.  This is the
established random-allocation/BIS mechanism
\citep{shenfeld2025randomallocation,dong2026lessrandom}.  We retain a
conservative direct bound with a separately structured high-precision
implementation; tighter privacy-loss distribution (PLD) accounting is prior
work \citep{feldman2026efficientallocation}.

\paragraph{Public-plan invariance.}
For Route P, $q,T,D$ do not derive from the observed retained-owner count, so
add/remove-adjacent inputs retain the same public mechanism.  For Route F,
$N,m,T,D$ define a fixed-population replace-one domain and are public constants,
not observed batch statistics.  For A, $t,k,E,D$ are public;
assignments and realized participation remain private.  The present research
profiles bind exact public benchmark fixtures; this restriction matters for
the conditional privacy boundary below.

\section{Evidence-Sealed Route-Label Consistency}

Route-label registration consistency is joint: a serialized contract can be
strict while its installed handler is internally mixed.
Table~\ref{tab:contract-surface} summarizes the shared contract surface.  We
therefore validate the handler before it may parse a contract.  ``Fixed'' below
means registered before protected execution, not estimated from realized
batches.

\begin{table}[t]
\centering
\small
\setlength{\tabcolsep}{3pt}
\begin{tabular}{@{}p{0.19\columnwidth}p{0.73\columnwidth}@{}}
\toprule
Layer & Required binding and rejection behavior \\
\midrule
Privacy & owner unit; exact adjacency; fixed $(\varepsilon,\delta)$ \\
Mechanism & sampler; $q$, $(N,m)$, or $(t,k,E)$; $T,C,\sigma,D$; sensitivity; empty-step rule \\
Data & single attribution; owner-local policy; exact mapping; public-fixed transform \\
Program & model, loss, optimizer, route executor, accountant, source bundle \\
Artifacts & typed public-plan digest; private mapping binding; model--run--summary--collection links \\
\bottomrule
\end{tabular}
\caption{Contract surface shared structurally but instantiated nominally by
each route.  An unsupported value or route mixture produces a named error.}
\label{tab:contract-surface}
\end{table}

\paragraph{Complete registration core.}
A handler $h$ contains route and schema identifiers, nominal tuple $\rho$,
exact contract type $\tau_c$ and compiled-plan type $\tau_P$, strict parser
$p$, compiler $c$,
route-specific registry object $r_h$, and witness map $W$.  Its complete core
is canonically serialized as
\[
 K(h)=\mathsf{canon}(\mathit{id},\mathit{schema},\rho,
 \tau_c,\tau_P,p,c,r_h,W).
\]
The seven required keys of $W$ are strict parser, strict compiler, accountant,
numeric cross-check, executor, postcompile integrity, and source binding; let
$\mathcal O$ denote exactly this key set.
Each witness names a callable, its repository path, and the exact file digest.
The evidence-report path and digest are outside $K(h)$ to avoid a circular
commitment; the handler binds them separately.

\paragraph{Registration judgment.}
Let $e$ be the route-local evidence report and write $\mathsf{SHA256}$ for
SHA-256.  We accept
$e\vdash h\ \mathsf{registered}$ only if
\begin{equation}
\begin{split}
\operatorname{dom}(W)&=\mathcal O,\quad\mathsf{Consistent}(h,W),\\
e.\mathit{status}&=\mathsf{PASS},\\
e.\mathit{id}&=h.\mathit{id},\quad
e.\mathit{schema}=h.\mathit{schema},\\
\mathsf{SHA256}(e.K)&=e.K_{\rm hash}=\mathsf{SHA256}(K(h)),
\end{split}
\label{eq:registration}
\end{equation}
and the separately bound report digest, report-payload digest, every live
source digest, and every callable/source resolution match.  Consistency also
requires the route-local registry's nominal fields and validated-executor
status, exact callable witnesses for parser/compiler/executor/integrity, and
distinct obligation keys.  This is a machine-checked registration predicate,
not a proof of privacy or arbitrary-code behavior.  Full route semantics and
obligations appear in Supplementary Material Secs.~2--3.

Figure~\ref{fig:v42-method-splice} summarizes the registered compilation path
and the handler-splice counterfactual evaluated below.
\begin{figure*}[t]
\centering
\includegraphics[width=0.9\textwidth]{figures/figure7.png}
\caption{Sealed compilation path and handler-splice counterfactual.  The top
lane shows exact-key handler selection, joint-core binding, branch-free generic
compilation/execution, and bound outputs.  The bottom lane reports accepted
nonendpoint splices exhaustively within the prespecified seven-surface binary
lattices; each of
P/F, P/A, and F/A has $2^7-2=126$ distinct hybrids, and all six unmodified
endpoint checks pass.  H is a separately registered external-accountant
stress test whose privacy-loss distribution implementation operates over A's
unchanged mechanism and executor; it is not a fourth mechanism.  Rejection
means only that a hybrid cannot retain either endpoint's registered label.}
\label{fig:v42-method-splice}
\end{figure*}

\paragraph{Generic parse and compile judgment.}
Let $R$ be an immutable map from exact route identifiers to registered
handlers, $x$ a serialized contract, $m$ a generated-record mapping, and $a$
a fixed preprocessor.  The generic judgment
\[
 R\vdash(x,m,a)\Downarrow(P_\rho,B_{\rm priv})
\]
first selects $h=R[x.\mathit{id}]$ and rechecks
Eq.~\eqref{eq:registration}.  It then requires exact schema, invokes $h.p$,
checks the exact returned contract type and nominal identity, invokes $h.c$,
checks the exact plan type and unchanged contract, and calls the registered
integrity method.  The dispatcher has no P/F/A identifier or type branch.
Route-specific compilation additionally checks the benchmark profile,
mapping/attribution, owner-local policy, fixed preprocessor, accounting, and
executor/source resolution.  The public plan $P_\rho$ omits realized owner
identities and sample counts; $B_{\rm priv}$ binds mappings and in-memory
preprocessor state.

\paragraph{Lemma 1 (handler-local conservative extension).}
Let $R'=R\uplus\{j\mapsto h_j\}$, where $j\notin\operatorname{dom}(R)$ and
$h_j$ passes registration.  If an old handler, its sources, and compilation
inputs are unchanged, then every terminating exact contract $x$ with
$x.\mathit{id}\in\operatorname{dom}(R)$ has the same parse and compilation
output under $R$ and $R'$.  Exact-key lookup selects the same immutable
handler, and every subsequent operation is handler-local and deterministic.
This syntactic property says nothing about privacy or the semantic correctness
of the newly added handler.  Evaluation checks exact objects and digests for
all six old P/F profile contracts.

\paragraph{Execution and artifact judgment.}
Execution is a second partial judgment,
\[
 (P_\rho,B_{\rm priv});\mathcal D \Downarrow
 (\theta_T,E_{\rm pub},E_{\rm priv}),
\]
which exists only if the compiled object passes before and after training, its
route loop completes, and the output verifier reconstructs artifact links.
Here $\theta_T$ is the final parameter vector, and $E_{\rm pub}$ and
$E_{\rm priv}$ are the public and private evidence artifacts.

Table~\ref{tab:judgment-evidence} summarizes where each lifecycle boundary is
enforced and what evidence persists.

\begin{table}[t]
\centering
\small
\setlength{\tabcolsep}{3pt}
\begin{tabular}{@{}p{0.17\columnwidth}p{0.32\columnwidth}p{0.41\columnwidth}@{}}
\toprule
Boundary & Enforced at & Persisted evidence \\
\midrule
registration & registry & core, report, and seven source witnesses \\
contract & parser/compiler & route, contract, and profile digests \\
data & compiler/executor & private map, policy, and transform bindings \\
accounting & compiler/verifier & two values and theorem/config identity \\
loop & executor/tests & declarations, diagnostics, and reference replay \\
links & serializer/verifier & model, execution, bundle, summary, collection \\
\bottomrule
\end{tabular}
\caption{Operational predicates and evidence boundaries.  Public artifacts
omit data-dependent identities and sampling diagnostics.}
\label{tab:judgment-evidence}
\end{table}

\paragraph{Strict data and executable binding.}
Structured contract files use duplicate-key detection, exact types, and no
extra fields.
Dataset contracts equal a versioned benchmark profile for every
profile-controlled field.  Compilation checks contiguous unique row indices,
unique generated-record identifiers, nonempty labels, exactly one owner per
record, and a full mapping digest.  The wrapper binds exact transformed
features, labels, and mapping before compilation and every run.  The affine
preprocessor checks dimensions and finite scales, declares no protected-data
fit, and is applied inside the executor.

Opacus already wraps core private-training objects
\citep{yousefpour2021opacus}.  Route P uses a custom owner executor and claims
only arithmetic compatibility: fixed-order R\'enyi differential privacy (RDP)
is recomputed through an Opacus-compatible calculation and
\texttt{dp-accounting} \citep{dpaccounting},
not executor equivalence.  Route F compares the Wang-style library result with
a direct implementation, using normalized noise $\sigma/2$
\citep{wang2019subsampled}.  Route A compares a floating-point direct
random-allocation calculation with a separately structured high-precision
implementation \citep{shenfeld2025randomallocation}.  The conservative
$\varepsilon$ must meet the target, the two route-internal values must agree
within the route tolerance, and the compiler binds the typed executor and
route-specific source bundle.  Agreement is implementation evidence; both
paths can share a mistaken theorem premise.

\paragraph{Execution, outputs, and threat model.}
Four domain-separated random streams cover initialization, owner sampling,
within-owner selection, and Gaussian noise.  The current backend uses NumPy's
PCG64 pseudorandom generator for reproducible research and records
\texttt{secure\_rng=false}.
Random coins, owner identities, mapping commitments, and realized sample
statistics remain private.  Public artifacts contain the fixed route,
accountant results, source manifest, public-benchmark metrics, and exact
float32 non-pickle model serialization.

The method targets accidental drift, stale results, misconfiguration, and
ordinary substitutions visible in checked inputs, source, or outputs.  It
assumes SHA-256 collision resistance and honest source, executor, evidence
builder, and verifier.  A malicious component can fabricate hashes, choose
weak noise, or exfiltrate private diagnostics.  Rejections and exceptions are
local controls, not DP outputs; exposing them can leak input-dependent
behavior.  Timing, hardware, and operating-system side channels are outside
scope.

\section{Conditional Privacy Scope}

\paragraph{Route-conditional statement.}
Assume a release procedure is total over its declared adjacent domain,
completes its fixed schedule without a data-dependent abort, keeps validation
outputs internal, fixes preprocessing and the owner-local rule independently
of protected input, contributes at most one vector of norm $C$ per
participating owner, and adds independent Gaussian noise at every step.  If P
uses its registered Bernoulli/add-remove accountant, F uses its registered
SRSWOR/replace-one bound with sensitivity $2C$, or A uses its registered
per-epoch uniform $k$-of-$t$ allocation/add-remove bound, then the applicable prior
analysis yields owner-level $(\varepsilon,\delta)$-DP for the final model under
that route's adjacency
\citep{wang2019subsampled,shenfeld2025randomallocation,
dong2026lessrandom}.  Fixed division, optimizer updates, and model release are
post-processing.  This restates the premises of existing analyses; the
registration judgment does not prove that an executor satisfies them.

\paragraph{What this statement does not establish.}
The present benchmark compilers require exact public-fixture mappings and can
reject changed inputs; they are not a total release algorithm over a private
adjacency domain.  All reported runs are therefore labeled
\texttt{research\_non\_release}, and accountant values are not release claims.
A deployment must replace PCG64 with audited secure randomness, make
compilation/failure behavior safe and total (or data-independently aborting),
and validate composition and the complete release procedure.  Hash equality
does not prove hidden random coins correct, and multiple releases require
composition.

\section{Evaluation}
\label{sec:evaluation}

We ask: (Q1) Do lifecycle bindings reject fixed failures beyond strong
contract-only parsing? (Q2) Does the complete-core seal reject splices accepted
by simpler registration boundaries? (Q3) Can it add a structurally different
route without changing old outputs, while separating controlled route and
contract surfaces? (Q4) Do route-specific accountants, executors, and artifact
checks support the registered endpoints?

\paragraph{Failure-oriented design.}
For Q1, contract-only means duplicate-key-safe loading plus exact parsing of
the registered contract; it excludes mapping/preprocessor compilation,
backend approval, data/execution binding, compiled integrity, and public
artifact redaction.  Each already-frozen case is assigned to its earliest
rejecting boundary.  For Q2, we compare three nested boundaries on the same
handler objects.  The weak boundary checks only a generic callable/type
interface.  The pre-seal boundary additionally checks nominal fields,
obligations, callable resolution, and source digests, but has no report
committing to the complete core.  The proposed boundary adds
Eq.~\eqref{eq:registration}.  Rejection means that a splice cannot retain
either endpoint's registered label; it does not declare the alternative
nonprivate.

\paragraph{Benchmarks and protocol.}
We use the UCI HAR dataset
\citep{anguita2013har,ucihar2013dataset}, the Wireless Sensor Data Mining (WISDM)
activity-recognition dataset
\citep{kwapisz2011activity}, and PhysioNet 2019 sepsis
\citep{physionet2019sepsis,reyna2020sepsis,pollard2026physionet,goldberger2000physionet}.  Their public owner-disjoint
protocols contain 21/9, 25/11, and 312/78 nominal train/test owners and
7,352/2,947, 7,049/3,331, and 1,601/398 generated train/test records.
Registered affine transforms produce 561, 24, and 240 features.  We use linear
classifiers and five fixed seeds per dataset and route.  All profiles target
$(\varepsilon,\delta)=(8,10^{-5})$; Table~\ref{tab:v38-profiles} gives the
nine public route profiles.  UCI HAR is a shared public substrate with the
concurrent study \citep{anonymous2026audit}: both studies
use its 7,352/2,947 rows, 561 features, 21/9 subject split, row order, labels,
and subject attribution.  Serialized mappings, route contracts, registered
sampling/accounting configurations, target pairs, models, results, and
artifact chains are separate.  Here UCI feeds owner-sampled P/F/A at
$(8,10^{-5})$; the concurrent work separately runs window-Poisson DP-SGD at
$(0.536827,10^{-6})$ and allows or blocks requested window/event/owner wording.
WISDM/sepsis mapping/evidence artifacts are likewise separate.  All results
reported here come from this paper's separately built pre-execution package.

\begin{table*}[t]
\centering
\small
\setlength{\tabcolsep}{4pt}
\begin{tabular}{@{}lllrrrr@{}}
\toprule
Dataset & Route & Adjacency / participation law & $T$ & $k$ or $m$ & $\sigma$ & $D$ \\
\midrule
UCI HAR & P & add/remove; Bernoulli $q=0.380952$ & 15 & -- & 1.269710 & 8 \\
 & F & replace-one; SRSWOR $m/N=8/21$ & 15 & $m=8$ & 3.806632 & 8 \\
 & A & add/remove; uniform $k$-of-$T$ allocation & 15 & $k=6$ & 1.229575 & 8 \\
\addlinespace[1pt]
WISDM & P & add/remove; Bernoulli $q=0.320000$ & 20 & -- & 1.251026 & 8 \\
 & F & replace-one; SRSWOR $m/N=8/25$ & 20 & $m=8$ & 3.838696 & 8 \\
 & A & add/remove; uniform $k$-of-$T$ allocation & 20 & $k=6$ & 1.090646 & 8 \\
\addlinespace[1pt]
Sepsis & P & add/remove; Bernoulli $q=0.100000$ & 50 & -- & 0.857344 & 32 \\
 & F & replace-one; SRSWOR $m/N=32/312$ & 50 & $m=32$ & 2.380381 & 32 \\
 & A & add/remove; uniform $k$-of-$T$ allocation & 50 & $k=5$ & 0.776173 & 32 \\
\bottomrule
\end{tabular}
\caption{Registered research profiles; all target
$(\varepsilon,\delta)=(8,10^{-5})$.  P samples each owner independently at
each step; F samples exactly $m$ distinct positions per step; A assigns each
owner independently to exactly $k$ distinct steps.  Every displayed A profile
uses one epoch, so $T=t$.  F uses replace-one
sensitivity $2C$; P and A use add/remove sensitivity $C$.  Route mechanisms
and accountants are prior work, not contributions of this paper.}
\label{tab:v38-profiles}
\end{table*}

\paragraph{Q1: strict-contract lifecycle counterfactual.}
The frozen P lineage contains 52 targeted mutations and eight reproduced
historical failures.  Contract-only rejects 36/52 and 5/8.  For the targeted
set, mapping/preprocessor compilation adds ten rejections, backend approval
adds one, and compiled integrity adds five.  For the historical set, exact
data/execution binding catches two stale mappings and the public-artifact
boundary catches one leaking certificate.  The full lifecycle therefore
reaches 52/52 and 8/8.  These retrospectively selected historical cases and the
fixed targeted corpus provide stage attribution, not held-out performance,
arbitrary-input recall, or a competitor benchmark.

\paragraph{Q2: discriminating registration ablation.}
We partition each handler into seven prespecified surfaces: identity/report,
contract and plan types, parser, compiler, route-local registry,
accountant/cross-check, and executor/integrity/source.  For each of P/F, P/A,
and F/A, all $2^7-2=126$ nonendpoint binary combinations are distinct.  As
Figure~\ref{fig:v42-method-splice} shows, weak validation accepts 378/378
and pre-seal validation accepts 42/378 (14 in every pair).  The complete-core
seal rejects 378/378, while all six endpoint checks across the three pairs pass.
A separately
implemented verifier that does not import the report builder reconstructs
every combination and decision.  Eight registration mutations, including
missing/duplicate obligations, changed source or report digests, spliced
parser/type, unvalidated executor, and duplicate route, also fail.
Thus the result is exhaustive conformance only over the declared
$3(2^7-2)=378$ author-specified nonendpoints, not a prevalence or open-world
sample.

\paragraph{Q3: route separation and bounded extension.}
We first control the scientific configuration rather than count
implementations.  In each P/F dataset pair, 38 fields, dataset and split,
privacy target, schedule,
clipping, contribution cap, model, optimizer, loss, and fixed preprocessor,
are exact matches.  Eight prespecified obligation families (route, schema,
adjacency, sampler, accountant/theorem, executor, noise, and the full
mechanism/accountant surface) are substituted in both directions over three
datasets.  All 48/48 nonvacuous substitutions fail: 12 at exact dispatch and
36 at the selected route's strict contract.  Thus the routes differ under a
matched shell, but this causal attribution covers only the 48 author-specified
substitutions and 38/38 fixed fields; this finite test is not semantic
completeness.

The generic parser/compiler contains none of the three route identifiers,
contract types, or an \texttt{if route\_id ==} branch; source inspection
confirms that it calls handlers.  All nine dataset--route contracts compile.
Adding A to the P/F registry leaves each of the six parsed contracts and
public plans structurally equal and preserves their exact digests, including
equality to the frozen predecessor outputs.

As a separate external-handler stress test, we register H, whose externally
maintained PLD implementation supplies a pessimistic bound for A's unchanged
mechanism and executor \citep{feldman2026efficientallocation}.  All four
handler endpoints validate, all nine old P/F/A contract--plan pairs remain
byte-exact,
and the complete-core seal rejects 756/756 nonendpoint hybrids over six
handler pairs.  This is one external accountant overlay, not arbitrary
third-party support.

We then partition every paired serialized contract into eight disjoint
surfaces: identity, execution, privacy/accounting, data lineage, schedule,
noise geometry, update contribution, and optimizer/sampler.  Across three
route pairs and three datasets, both endpoints parse and all
$9(2^8-2)=2{,}286$ nonendpoints fail.  Strict route parsers explain much of
this result, so it supports route separation rather than the central novelty.
The older P add/remove test also holds $q,T,D$ and both accountant values fixed
while changing only the private mapping commitment.
Full counterfactual constructions and tables appear in Supplementary Material
Sec.~4.

\paragraph{Q4: accountant, execution, and artifact correspondence.}
For A, the production direct calculation and a separately structured
high-precision implementation agree for all three profiles, both add/remove
privacy directions, and every registered order; six mutations to the bound
accountant-query identity and thirteen invalid inputs are exercised.  In the
frozen A gate, external PLD
values are recorded only as non-authorizing comparisons; H is the separately
registered external-accountant stress test.  A reference executor that imports
no production implementation then matches 15/15 runs, 425/425 steps,
9,180/9,180 owner vectors, and 30/30
final tensors.  A separate fixture contains a true empty step and confirms
that Gaussian noise is generated for both model tensors and the optimizer
advances.  P/F
retain their 30/30 exact run replays.  The A public archive rebuilds
deterministically and contains no checked private assignment/identity fields.
Together the 15 A and 30 P/F runs complete a 45/45 fixed-seed integration grid.
It tests compilation/execution correspondence, not an outcome distribution;
different adjacency and calibration preclude utility or inferential comparison.
Accountant and executor details appear in Supplementary Material Sec.~6.

\paragraph{Evidence authority.}
Lifecycle counts are retrospective attribution; lattice counts are exhaustive
only within named author-specified finite universes; 48/48 is matched-shell
ablation; and 45/45 is fixed-seed integration, not a utility or population
sample.  Full-universe deterministic predicates have no sampling error, so
significance tests do not apply.  Table~\ref{tab:v38-conformance} binds every
result to an artifact, verifier, and exact denominator.

\begin{table*}[t]
\centering
\small
\setlength{\tabcolsep}{5pt}
\begin{tabular}{@{}p{0.68\textwidth}r@{}}
\toprule
Evidence role, bound artifact; verifier/comparator & Exact result \\
\midrule
Endpoint conformance, route reports; registration verifier &
3/3; 7/7 per handler \\
Lifecycle attribution, frozen P manifest; stage verifier &
36/52 vs. 52/52; 5/8 vs. 8/8 \\
Matched-shell ablation, P/F substitution report; separate verifier &
48/48; 38/38 fixed \\
Compilation/extension, contracts, plans, and frozen P/F/H outputs;
compiler/digest verifiers & 9/9; 6/6; H 4/4, 9/9 \\
Finite-lattice conformance, three-, six-, and nine-pair reports; separate
verifier & 378/378; 756/756; 2,286/2,286 \\
A integration, trace bundle and empty-step fixture; reference
executor/comparator & 15/15; 425/425; 9,180/9,180; 30/30; 1/1 \\
P/F integration, trace bundles; reference executor &
30/30; 850/850; 60/60 \\
Regression conformance, focused/full logs; test runner &
44/44; 154/154 run; 7/161 skip \\
\bottomrule
\end{tabular}
\caption{Artifacts, verifiers, and exact denominators.  ``Exhaustive'' means
only the named lattice; 45 fixed-seed runs test integration, not utility or a
population.  No result proves privacy, arbitrary-input refinement, or
adversarial attestation.}
\label{tab:v38-conformance}
\end{table*}

\paragraph{Reproducibility controls.}
Table generators reject duplicate serialization keys and recompute displayed
aggregates.  Separate verifiers reclassify all 60 contract-only cases,
reconstruct the 48 matched-shell substitutions, handler and contract lattices,
conservative-extension rows, route evidence, private diagnostics, source
links, and registered executor traces.  Focused route-accounting, compilation,
execution, registration, and lattice tests pass 44/44; the full frozen
regression collects 161 tests: all 154/154 executed tests pass, and 7/161
raw-data-scoped cases skip.
Two optional pandas dependency-version warnings are separate from those skips.

\section{Related Work}

Table~\ref{tab:v40-ownership} summarizes the closest-work boundary;
Supplementary Material Sec.~5 gives the citation-level comparison.

\begin{table*}[t]
\centering
\small
\setlength{\tabcolsep}{4pt}
\begin{tabular}{@{}p{0.18\textwidth}p{0.78\textwidth}@{}}
\toprule
Prior-work family & Established scope; boundary evaluated here \\
\midrule
Sampling/accounting &
Units and P/F/A laws/accounting; here: finite registration across established
endpoints. \\
Formal/typed DP &
Semantic types, proofs, and refutations; here: finite route registration
without refinement. \\
Testing/verification &
Falsification and plan equivalence; here: pre-execution splice conformance. \\
Provenance/builds &
Signed layouts and reproducible builds; here: route-label commitment and the
378/42/0 ablation. \\
Attestation/integrity &
Adversarial execution evidence; here: an honest-source registration seal. \\
\bottomrule
\end{tabular}
\caption{Closest-work scope and evaluated boundary: complete-core registration and
the 42-to-0 handler counterfactual.}
\label{tab:v40-ownership}
\end{table*}

\paragraph{Privacy units, sampling, and accounting.}
User-level grouping and bounding are prior \citep{chua2024mind}; so are RDP
composition and P/F sampling analyses
\citep{mironov2017renyi,wang2019subsampled,birrell2024fixed}.  Random allocation,
its PLD accounting, and BIS's exactly-$k$ participation law are also prior
\citep{shenfeld2025randomallocation,feldman2026efficientallocation,
dong2026lessrandom}.  Add/remove and replace-one adjacency protect different
facts \citep{pradhan2026beyond}, while multi-attribution needs another analysis
\citep{ganesh2025multiattribution}.  We therefore use P/F/A only as
heterogeneous endpoints, claiming no mechanism or accountant novelty.

\paragraph{Formal DP frameworks and implementation guidance.}
Formal analyzers provide DP types, proofs, or refutations
\citep{zhang2017lightdp,wang2020checkdp,near2019duetfull,
haselwarter2026modularfull}; OpenDP supplies stronger domains, metrics, and
stability/privacy maps \citep{hay2020opendp}.  Modular implementation review is
prior \citep{kifer2020guidelines}, as is P's Opacus stack
\citep{yousefpour2021opacus}.  We claim none of those semantics.  Our finite
object instead binds the generated-record map, application profile, handlers,
source, and lifecycle, then measures splice rejection and old-route
preservation without a refinement proof.

\paragraph{Software provenance and reproducible builds.}
in-toto authenticates signed layouts, links, and artifact flow
\citep{torresarias2019intoto}; reproducible builds seek bit-identical rebuilds
\citep{lamb2022reproducible}.  Either can strengthen our chain, but neither alone
rejects a mixed route.  Our residual is the DP route-label predicate and
378/42/0 counterfactual, not authentication, generic provenance, or
source-to-binary reproducibility.

\paragraph{Testing, training verification, and attestation.}
Sampler/accountant mismatches can materially weaken implementations
\citep{chua2024private}.  DPCheck, DP-Auditorium, and Re:cord-play test or
falsify implementations
\citep{zhang2020dpcheck,kong2024dpauditorium,cebere2026recordplay}; TrainCheck
detects runtime violations and TrainVerify proves plan equivalence
\citep{jiang2025training,lu2025trainverify}.  VeriDP and the integrity SoK cover
attestation and broader evidence models
\citep{abdolmaleki2026veridp,rizzini2026integrity}.  We instead evaluate finite
pre-execution registration and extension, claiming no detection, proof, or
attestation.

\paragraph{Concurrent phase boundary.}
The concurrent validator decides wording from frozen runs
\citep{anonymous2026audit}.  We instead consume a pre-execution contract and
emit a plan, model, and run chain; neither pipeline consumes the other's
outputs, and we emit no post-hoc report.

\section{Limitations}

Our evidence supports a finite, author-constructed, honest-execution
registration boundary, not a general DP compiler.  The 378/756 handler
splices and 2,286 contract nonendpoints exhaust only their prespecified
universes; the 60-case P lifecycle and 48 matched-shell substitutions are fixed
controls, not arbitrary-input recall estimates.  Unmodeled obligations, new
adjacencies, mechanisms, models, or handlers require new analyses, executors,
and evidence.  Within this scope, the controls separate the roles of strict
parsing, source binding, and complete-core sealing.

SHA-256 commits declared artifacts but attests neither runtime behavior nor
random coins; deployment needs secure randomness, safe failure handling, and
stronger provenance or independent-rebuild guarantees.  Linear-model, five-seed
replays test compilation/execution correspondence, not scalability, leading
utility, or comparisons among P/F/A laws.  H tests one external accountant over
A; its environment is bound but neither redistributed nor rerun.  Full profiles,
evidence identifiers, commands, and further limitations are in Supplementary
Material Secs.~7--10.

\section{Conclusion}

Complete-core registration is one route-installation judgment.  Lifecycle checks
catch 16 targeted and three historical cases; the matched shell rejects 48
substitutions.  The seal closes 42 splices, rejects 378 nonendpoints, and preserves
six outputs after adding A.  Execution cannot establish owner-level DP; these
tests prove neither privacy nor a general verifier.  \textbf{AI disclosure:}
Generative-AI tools assisted editing; authors verified all content and remain
responsible.

\bibliography{references_v28_candidate,references_v30_additions,references_v33_additions,references_v35_additions,references_v37_additions}

\end{document}
